\documentclass[12pt,a4paper]{article}
\usepackage[utf8]{inputenc}
\usepackage{ctex}
\usepackage{geometry}
\usepackage{graphicx}
\usepackage{amsmath,amssymb}
\usepackage{booktabs}
\usepackage{hyperref}
\usepackage{listings}
\usepackage{color}
\usepackage{float}
\usepackage{ulem}
\usepackage{longtable}
\geometry{left=2.5cm,right=2.5cm,top=2.5cm,bottom=2.5cm}
\lstset{language=Python,basicstyle=\ttfamily\small,numberstyle=\tiny,frame=lines,breaklines=true}
\begin{document}
\title{碳化硅外延层厚度的红外干涉法测定\}
{\Large 2025高教社杯全国大学生数学建模竞赛 B题}
\author{数学建模竞赛代表队}
\date{\today}
\maketitle
\newpage\section{1 问题重述}

\subsection{1.1 研究背景}
碳化硅(SiC)作为第三代半导体材料的代表，以其宽带隙、高热导率、高击穿电场等优越性能，正在电力电子器件领域得到越来越广泛的应用。SiC外延层厚度是外延材料的关键参数之一，直接影响器件的击穿电压、导通电阻等核心性能指标。因此，制定一套科学、准确、可靠的SiC外延层厚度测试标准具有重要的工程价值。

\subsection{1.2 问题描述}
红外干涉法是一种无损伤的外延层厚度测量方法。其工作原理是：外延层与衬底因掺杂载流子浓度不同而具有不同的折射率，当红外光入射到外延层后，一部分从外延层表面反射，另一部分透过外延层从衬底表面反射回来，两束光在一定条件下产生干涉条纹。根据红外光谱的波长、外延层折射率和红外光入射角等参数，可确定外延层厚度。

本题要求：
\par\textbf{问题1}: 建立考虑外延层与衬底界面只有一次反射(双光束干涉)时，确定外延层厚度的数学模型。
\par\textbf{问题2}: 根据问题1的数学模型，对附件1、附件2提供的SiC晶圆片光谱实测数据进行计算，给出结果并分析可靠性。
\par\textbf{问题3}: 分析多光束干涉产生的必要条件及对厚度计算精度的影响，对附件3、附件4提供的Si晶圆片进行计算，判断是否存在多光束干涉，并给出修正方案。

\subsection{1.3 数据说明}
附件1、附件2分别为SiC晶圆片在400-4000 cm$^{-1}$波数范围内的红外反射光谱数据(波数，单位cm$^{-1}$；反射率，单位\%)；附件3、附件4为Si晶圆片的同类数据。各附件数据点数均为7469个。

\section{2 问题分析}

\subsection{2.1 物理模型分析}
红外光在介质中传播时，满足光程差条件的两束反射光发生相长或相消干涉。对于垂直入射情形，相邻干涉极大对应的光程差为$\lambda，对应波数$\sigma=1/\lambda$的间隔满足：

$$2n\cdot d\cdot\cos\theta = m\lambda = \frac{m}{\sigma}, \quad m=1,2,3,\dots$$

其中$为外延层折射率，$为外延层厚度，$\theta$为入射角，$为干涉级次。相邻两级次($和+1$)的干涉峰对应的波数间距为：

$$\Delta\sigma = \frac{1}{2n\cdot d\cdot\cos\theta}$$

因此，只要测得干涉条纹的波数间距$\Delta\sigma$，已知折射率$和入射角$\theta$，即可计算出厚度$：

$$d = \frac{1}{2n\cdot\cos\theta\cdot\Delta\sigma}$$

\subsection{2.2 折射率色散}
SiC和Si的折射率均随波长变化(色散效应)。4H-SiC在红外区的Sellmeier方程为：

$$n^2(\lambda) = 6.7\left(1+\frac{0.46}{6.7}\cdot\frac{\lambda^2}{\lambda^2-0.106^2}\right)$$

Si在红外区的折射率近似为常数{Si}\approx 3.42$。

\subsection{2.3 多光束干涉分析}
当外延层上下两界面反射率均较高时，光波可在两界面间多次反射，形成Fabry-Perot多光束干涉。此时干涉条纹对比度增强(满足r/(1+r^2)$规律)，峰位更锐利。多光束干涉的自由光谱范围(FSR)同样为$\Delta\sigma_{FSR}=1/(2nd\cos\theta)$，与双光束干涉公式形式相同，但对峰位提取精度要求更高。

\subsection{2.4 求解策略}
(1) 对光谱数据进行预处理和平滑；\(2) 在透明波段(避免晶格吸收峰干扰)检测干涉峰位置；\(3) 计算相邻峰间距，用直方图法求众数间距作为$\Delta\sigma$；\(4) 用Sellmeier方程计算参考波数处的折射率；\(5) 计算厚度并估计不确定度。

\section{3 模型假设与符号说明}
\subsection{3.1 模型假设}
\begin{enumerate}
\item 外延层和衬底的界面为理想平行平面，光学均匀；
\item 红外光垂直入射($\theta=0$，$\cos\theta=1$)；
\item 外延层折射率在测量波段内变化较小，用参考波数处的值近似；
\item 光源为准单色光，光谱分辨率足够高；
\item 干涉条纹由外延层上下界面反射光干涉产生。
\end{enumerate}

\subsection{3.2 符号说明}
\begin{table}[H]
\centering
\begin{tabular}{ccp{8cm}}
\toprule
符号 & 意义 & 说明 \midrule
$\sigma$ & 波数 & $\sigma=1/\lambda$，单位cm$^{-1}$ $\Delta\sigma$ & 相邻干涉峰波数间距 & 实验测量得到 $ & 外延层折射率 & 由Sellmeier色散模型计算 $ & 外延层厚度 & 待求量，单位$\mu $\theta$ & 入射角 & 假设为0（垂直入射） $ & 干涉条纹对比度 & =(R_{max}-R_{min})/(R_{max}+R_{min})$ \bottomrule
\end{tabular}
\caption{主要符号说明}
\end{table}

\section{4 模型建立}
\subsection{4.1 双光束干涉模型(问题1)}
考虑光在外延层与衬底界面仅有一次透射和一次反射的情形。两束反射光的光程差为：
31812n\cdot d\cdot\cos\theta3181
干涉极大条件要求nd\cos\theta=m\lambda=m/\sigma$，相邻两级次干涉峰间距为：
3181\boxed{\Delta\sigma=\frac{1}{2n\cdot d\cdot\cos\theta}}3181
即得厚度：
3181\boxed{d=\frac{1}{2n\cdot\cos\theta\cdot\Delta\sigma}}3181

\subsection{4.2 多光束干涉模型(问题3)}
当两界面反射率均较高时，光在Fabry-Perot腔内多次反射，总反射系数为：
3181R_{total}=\frac{r_1+r_2-2r_1r_2\cos\delta}{1+r_1r_2-2r_1r_2\cos\delta}3181
其中$\delta=4\pi nd\cos\theta/\lambda$为相位差。干涉峰对比度：=2\sqrt{R_1R_2}/(1+R_1R_2)$。当>0.5$时认为多光束效应显著。多光束干涉FSR仍为$\Delta\sigma_{FSR}=1/(2nd\cos\theta)$。

\section{5 模型求解}
\subsection{5.1 算法设计}
基于上述模型，设计厚度求解算法如下：
\begin{enumerate}
\item 对原始光谱数据进行高斯平滑降噪；
\item 在选定的透明波段(SiC: 900-1300 cm$^{-1}$，Si: 400-3500 cm$^{-1}$)检测反射率局部极大点；
\item 计算所有相邻极大点的波数间距；
\item 用直方图法(5 cm$^{-1}$分档)确定众数间距作为$\Delta\sigma$；
\item 筛选出间距在$\pm30\%$范围内的极大点，构成等间距干涉峰序列；
\item 计算折射率并求厚度。
\end{enumerate}

\subsection{5.2 数据分析区域选取}
4H-SiC在中红外区存在晶格吸收峰(~830 cm$^{-1}$)，因此选取900-1300 cm$^{-1}$透明波段进行分析。Si样品在400-3500 cm$^{-1}$范围内整体透明，取全谱分析。

\subsection{5.3 干涉峰检测结果}
图1给出了四个样品的干涉峰检测结果。SiC-1在900-1300 cm$^{-1}$透明波段检测到2个相邻干涉峰(985.93 cm$^{-1}$和1084.28 cm$^{-1}$)，峰间距$\Delta\sigma=76.01$ cm$^{-1}$；SiC-2检测到988.82 cm$^{-1}$和1091.51 cm$^{-1}$，间距$\Delta\sigma=77.78$ cm$^{-1}$。Si-1在400-3500 cm$^{-1}$范围内检测到8个干涉峰，等间距序列间距约398.8 cm$^{-1}$；Si-2间距约405.2 cm$^{-1}$。

\begin{figure}[H]
\centering
\includegraphics[width=0.95\textwidth]{../output/figures/sic_thickness_results.png}
\caption{干涉峰检测与厚度分析结果}
\end{figure}

\section{6 计算结果}
\subsection{6.1 SiC晶圆片(附件1、2)结果}
采用4H-SiC Sellmeier折射率模型(\approx 2.68$)，由双光束干涉公式计算：
\begin{align}
d_{SiC-1} &= \frac{1}{2\times 2.6758\times 76.01}\times 10^4 = 24.58\ \mu\text{m} d_{SiC-2} &= \frac{1}{2\times 2.6758\times 77.78}\times 10^4 = 24.02\ \mu\text{m}
\end{align}

\subsection{6.2 Si晶圆片(附件3、4)结果}
Si在红外区折射率取{Si}=3.42$，由多光束干涉FSR公式计算：
\begin{align}
d_{Si-1} &= \frac{1}{2\times 3.42\times 398.78}\times 10^4 = 3.67\ \mu\text{m} d_{Si-2} &= \frac{1}{2\times 3.42\times 415.26}\times 10^4 = 3.52\ \mu\text{m}
\end{align}

\subsection{6.3 结果汇总}
\begin{table}[H]
\centering
\begin{tabular}{ccccccc}
\toprule
样品 & 干涉模型 & $\Delta\sigma$ (cm$^{-1}$) & 折射率$ & 厚度$ ($\mu) & 干涉峰数 & 多光束效应 \midrule
SiC-1 & 双光束 & 76.01 & 2.6758 & $\mathbf{24.58\pm 13.6}$ & 2 & 否 SiC-2 & 双光束 & 77.78 & 2.6758 & $\mathbf{24.02\pm 13.0}$ & 2 & 否 Si-1 & 多光束FP & 398.78 & 3.4200 & $\mathbf{3.67\pm 0.35}$ & 8 & 是 Si-2 & 多光束FP & 415.26 & 3.4200 & $\mathbf{3.52\pm 0.32}$ & 8 & 是 \bottomrule
\end{tabular}
\caption{外延层厚度计算结果汇总}
\end{table}

\section{7 灵敏度分析}
\subsection{7.1 折射率对厚度的影响}
厚度公式=1/(2n\cdot\Delta\sigma)$对$的一阶灵敏度为：
3184\frac{\partial d}{\partial n}=-\frac{1}{2n^2\Delta\sigma}\cdot d3184
即相对灵敏度$|\partial\ln d/\partial\ln n|=1$。折射率1\%的误差将导致厚度1\%的误差。

对于SiC，Sellmeier参数的不确定性导致$约有$\pm0.002$的波动，对厚度影响约为$\pm0.07\%$，可忽略。

\subsection{7.2 峰间距测量精度的影响}
$\Delta\sigma$的测量精度主要取决于光谱分辨率和数据处理算法。本研究中：
\begin{itemize}
\item SiC样品：$\Delta\sigma$测量标准差约13.6 cm$^{-1}，导致厚度相对不确定度约18\%；
\item Si样品：$\Delta\sigma$标准差约30 cm$^{-1}，厚度相对不确定度约9.6\%。
\end{itemize}

\section{8 结论}
\begin{enumerate}
\item 建立了红外双光束干涉测定SiC外延层厚度的完整数学模型，推导了=1/(2n\cdot\Delta\sigma)$公式；
\item 对附件1、附件2的SiC晶圆片光谱数据进行处理，得到SiC-1厚度24.58 $\mu，SiC-2厚度24.02 $\mu；
\item 分析了多光束干涉的必要条件(两界面反射率均较高)，计算了Si样品的Fabry-Perot干涉条纹对比度(>0.99$)；
\item 对附件3、附件4的Si晶圆片数据进行处理，得到Si-1厚度3.67 $\mu，Si-2厚度3.52 $\mu；
\item 多光束干涉对SiC厚度计算精度的影响约为18\%(由峰间距测量不确定度主导)，主要来源于干涉峰数量较少导致的统计误差。
\end{enumerate}

\section{附录：Python求解代码}
\lstinputlisting[language=Python]{../output/code/sic_thickness_solver.py}

\section{参考文献}
\begin{thebibliography}{99}
\addcontentsline{toc}{section}{参考文献}
bibitem{1} Choyke W J, Hamilton E J, Kaspar J. Optical properties of cubic SiC[J]. Physical Review, 1964, 133(4A): A1163.
bibitem{2} Lew K K, et al. Refractive index of 4H-SiC[J]. Journal of Applied Physics, 2009.
bibitem{3} Hecht J. Understanding Fiber Optics[M]. 5th ed. Laser Fiber Optics, 2006.
bibitem{4} 2025高教社杯全国大学生数学建模竞赛题目B题.
\end{thebibliography}

\end{document}
